% =============================================================================
%  KẾT LUẬN
% =============================================================================
\chapter*{KẾT LUẬN}
\addcontentsline{toc}{chapter}{KẾT LUẬN}
\markboth{KẾT LUẬN}{}

\section*{1. Tóm tắt kết quả nghiên cứu}
\addcontentsline{toc}{section}{1. Tóm tắt kết quả nghiên cứu}

Luận văn đã giải quyết bài toán tối ưu phân bổ tài nguyên đường xuống
trong mạng vô tuyến STAR-RIS hỗ trợ RSMA, một trong những bài toán
trung tâm cho mạng vô tuyến 6G trong tương lai. Nội dung chính của
Luận văn được tóm tắt như sau:

\begin{enumerate}[label=\textbf{(\arabic*)},leftmargin=*]
\item \textbf{Mô hình MISO STAR-RIS--RSMA.} Hệ thống sử dụng BS $M=4$
anten, beamforming phức cho luồng chung và các luồng riêng, STAR-RIS ES và
kênh động Gauss--Markov.
\item \textbf{Bài toán P0 và surrogate có ràng buộc kỳ vọng.} QoS được xử
lý bằng biến đối ngẫu per-user, augmented penalty và lịch trình two-stage
dual freeze; metric per-user và all-users được báo cáo tách biệt.
\item \textbf{MADDPG-CTDE ba tác tử.} Actor sử dụng quan sát cục bộ kích
thước $73$, $577$, $401$; Critic nhận trạng thái canonical $681$ chiều và
joint action $140$ chiều, không lặp các đặc trưng dùng chung.
\item \textbf{Tính đúng và tái lập.} Quy ước kênh và prior pha được kiểm
chứng bằng oracle tín hiệu thu; PPO xử lý đúng terminal/truncation; replay,
checkpoint, seed split và source hash được khóa bằng test tự động.
\item \textbf{Trạng thái kết quả.} Các bảng và con số từ pipeline SISO/tĩnh
cũ chỉ được giữ để audit. Kết luận định lượng cuối cùng chỉ được cập nhật sau
khi chạy đủ ma trận MISO trên commit freeze và hoàn tất paired tests với hiệu
chỉnh Holm. Không tuyên bố MADDPG vượt TD3 hoặc TD3-Matched trước bước này.
\end{enumerate}

\end{enumerate}

\section*{2. Đóng góp của Luận văn}
\addcontentsline{toc}{section}{2. Đóng góp của Luận văn}

\noindent\textbf{Về mặt lý luận:}
\begin{itemize}[leftmargin=*]
\item Mở rộng bài toán P0 cổ điển với ràng buộc QoS trên từng người dùng
và biên độ STAR-RIS như biến tối ưu, phản ánh chính xác hơn yêu cầu
thực tế của 6G.
\item Cung cấp baseline căn pha giải tích nhất quán với phương trình tín
hiệu thu MISO; biến thể residual được giữ như ablation tùy chọn và không
được đồng nhất với cấu hình kết quả chính.
\item Đề xuất hàm phần thưởng Lagrangian surrogate với các biến đối
ngẫu per-user cập nhật bằng projected dual gradient, một thiết kế có
thể áp dụng cho các bài toán DRL có ràng buộc kỳ vọng.
\item Đề xuất lịch trình heuristic hai pha (two-stage dual freeze)
làm giảm tính non-stationary của phần thưởng ở giai đoạn cuối huấn
luyện, giúp đường cong huấn luyện dễ diễn giải; đây là một heuristic
ổn định hóa, không phải một cơ chế có bảo đảm hội tụ điểm yên ngựa.
\item Cung cấp bằng chứng thực nghiệm có kiểm định thống kê nghiêm ngặt
(paired test $+$ Holm, $8$ hạt giống/điểm) \emph{kiểm chứng} giả thuyết
``lợi thế phân rã CTDE tăng theo kích thước không gian hành động'':
quét $N$ từ $16$ đến $128$ cho thấy giả thuyết này \emph{không} được
ủng hộ --- MADDPG tương đương TD3 và TD3-Matched ở mọi $N$. Kết quả
âm tính được báo cáo trung thực này bản thân nó là một đóng góp về
phương pháp đánh giá công bằng (đơn vị mẫu mức hạt giống, baseline
khớp tham số, hiệu chỉnh đa so sánh).
\end{itemize}

\noindent\textbf{Về mặt thực tiễn:}
\begin{itemize}[leftmargin=*]
\item Cung cấp một thư viện mã nguồn mở (Python/PyTorch) đầy đủ, có
khả năng tái lập, sẵn sàng để mở rộng cho các nghiên cứu tiếp theo.
\item Cấu hình mô phỏng $N=32, K=4, P_{\max}=30$ dBm trở thành một
chuẩn so sánh tham chiếu hữu ích cho cộng đồng nghiên cứu STAR-RIS-RSMA.
\item Chứng minh khả năng triển khai DRL cho phân bổ tài nguyên 6G
trong điều kiện thời gian thực (độ trễ cỡ mili-giây, nhanh hơn $13$
lần so với heuristic AO-Grid giải lại cho từng thể hiện kênh; số liệu
chính thức dùng cặp bảng CPU/GPU đo riêng).
\end{itemize}

\section*{3. Hạn chế của nghiên cứu}
\addcontentsline{toc}{section}{3. Hạn chế của nghiên cứu}

Để đánh giá khách quan đóng góp của Luận văn, phần này thảo luận chi
tiết các hạn chế và giới hạn áp dụng của khuôn khổ đề xuất, đồng thời
chỉ ra các điều kiện biên mà ngoài đó các kết quả không còn được bảo
đảm.

\textbf{Giới hạn quy mô MISO.} Cấu hình chính sử dụng $M=4$ anten,
không đại diện cho massive MIMO. Khi $M$ tăng, số biến beamforming, kích
thước trạng thái và chi phí huấn luyện tăng mạnh; khả năng mở rộng theo $M$
phải được đánh giá riêng và không được suy diễn từ sweep theo $N$.

\textbf{Giả thiết CSI hoàn hảo.} Giả thiết A1 (Bảng \ref{tab:assumptions})
là tiêu chuẩn trong các nghiên cứu mốc so sánh hiệu năng, nhưng không
phản ánh thực tế. Việc ước lượng kênh $\mathbf{h}_{r,k}^X$ giữa STAR-RIS và
UE đặc biệt khó khăn vì STAR-RIS thụ động không có RF chain để xử lý
pilot. Các kỹ thuật ước lượng kênh chuyên biệt cho RIS (như cascaded
channel estimation) có sai số điển hình $5$--$10\%$
\cite{wu2024deep}, có thể ảnh hưởng đáng kể đến hiệu năng MADDPG trong
thực tế \cite{cernaloli2023meta}.

\textbf{Mô hình kênh đơn giản hóa.} Pha-đinh Rayleigh i.i.d.\ là mô
hình tốt cho môi trường tán xạ phong phú nhưng không phản ánh các kịch
bản 6G phức tạp hơn như near-field communication, mmWave/THz với
clustering effect, hay correlated fading. Mô hình kênh dựa trên hình
học (geometry-based stochastic model) sẽ phù hợp hơn cho các tình
huống này, nhưng đòi hỏi mở rộng tham số hệ thống.

\textbf{Cấu hình một cell.} Luận văn chỉ xét một cell đơn lẻ với một
BS và một STAR-RIS. Trong thực tế, các mạng đa cell có hiện tượng nhiễu
liên cell (inter-cell interference), handover, và phối hợp đa BS
(CoMP). Mở rộng sang đa cell yêu cầu cơ chế phối hợp giữa các MADDPG
controller, có thể được nghiên cứu thông qua federated MARL.

\textbf{Kiến trúc mạng gắn với cấu hình huấn luyện.} Mạng nơ-ron
Actor và Critic có kích thước cố định theo cấu hình $(N, K)$; khi
$N$ hoặc $K$ thay đổi, cần huấn luyện lại từ đầu (thí nghiệm quét
$N$ ở Chương \ref{ch:chuong4} thực hiện đúng theo cách này). Kiến
trúc graph neural network hoặc transformer-based policy có thể học
chính sách tổng quát hóa theo kích thước, nhưng chưa được khảo sát
trong Luận văn này.

\textbf{Hạn chế phần cứng STAR-RIS.} Pha STAR-RIS được giả định liên
tục trong $[0, 2\pi)$. Phần cứng STAR-RIS hiện tại chỉ hỗ trợ pha
lượng tử với $1$--$3$ bit (tương ứng $2, 4, 8$ mức pha rời rạc). Tác
động lượng tử hóa lên hiệu năng MADDPG chưa được khảo sát; có thể yêu
cầu post-processing rời rạc hóa hoặc huấn luyện với quantized action
space.

\textbf{Đánh giá chỉ qua mô phỏng số.} Tất cả kết quả được thu nhận
qua mô phỏng số trên Python/PyTorch, chưa có hardware-in-the-loop
validation hay SDR testbed evaluation. Việc kiểm chứng trên phần cứng
thực sẽ phát hiện các yếu tố không được mô hình hóa như jitter,
clock drift, hardware impairment ở BS và STAR-RIS.

\textbf{Cỡ mẫu và thiết kế thống kê.} Cả cấu hình chuẩn lẫn thí nghiệm
quét $N$ đều dùng $8$ hạt giống huấn luyện/điểm --- tiệm cận khuyến
nghị $N \geq 10$ của cộng đồng DRL; các kết luận chính dùng thống kê
theo cặp (paired test) trên các trung bình mức hạt giống với hiệu
chỉnh Holm--Bonferroni. Với cỡ mẫu $n=8$, lực kiểm định còn hạn chế:
hiệu số MADDPG--DDPG tuy có kích thước hiệu ứng lớn ($d=0{,}98$) vẫn
chưa đạt ngưỡng ý nghĩa sau hiệu chỉnh ($p_{\text{Holm}}=0{,}081$), nên
một số so sánh cần cỡ mẫu lớn hơn để kết luận chắc chắn. Đáng chú ý,
baseline TD3-Matched (tác tử đơn có tổng tham số khớp MADDPG) cho phép
tách bạch vai trò của \emph{phân rã đa tác tử} khỏi dung lượng mạng:
kết quả cho thấy MADDPG không vượt TD3-Matched có ý nghĩa, nên lợi thế
(nếu có) của phân rã không thể được khẳng định trong khảo sát này.

\section*{4. Hướng nghiên cứu tiếp theo}
\addcontentsline{toc}{section}{4. Hướng nghiên cứu tiếp theo}

Dựa trên phân tích các hạn chế, Luận văn đề xuất bảy hướng nghiên cứu
tiếp theo, được sắp xếp theo mức độ khả thi và ưu tiên thực tiễn.

\textbf{H1. Khuôn khổ MADDPG cho MIMO BS với beamforming.} Mở rộng tác
tử BS Power Agent thành tác tử Joint Power-Beamforming Agent, xuất ra
cả $\{P_c, P_k, c_k\}$ và các vector tiền mã hóa $\mathbf{w}_k \in
\mathbb{C}^{M \times 1}$. Kích thước hành động tăng tuyến tính theo
$M$, nhưng cấu trúc phân rã ba tác tử vẫn duy trì được. Hướng này
giúp tận dụng các băng tần mmWave/THz với mảng anten lớn
\cite{wu2024deep}, \cite{wang2024rsbnn}.

\textbf{H2. Xử lý CSI không hoàn hảo bằng Robust RL hoặc Meta-RL.}
Hai hướng tiếp cận chính: (i) huấn luyện MADDPG với noise injection
trên CSI để học chính sách bền vững (Robust RL); (ii) sử dụng
meta-learning để học chính sách có khả năng thích nghi nhanh với phân
phối CSI mới sau vài bước cập nhật \cite{faramarzi2024meta},
\cite{cernaloli2023meta}. Hướng (ii) đặc biệt phù hợp với mạng vô
tuyến không ổn định.

\textbf{H3. MARL phân tán hoàn toàn (Fully Distributed MARL).} Hiện
tại các Actor vẫn chia sẻ replay buffer trong huấn luyện. Trong các
kịch bản đa cell hoặc đa nhóm tác tử, mỗi tác tử cần có replay buffer
cục bộ và chỉ truyền thông giới hạn với các tác tử khác. Cách tiếp cận
này giảm độ trễ truyền thông và tăng tính riêng tư \cite{mirza2024multi}.

\textbf{H4. Hybrid AO--MADDPG initialization.} Sử dụng nghiệm AO local
search (hoặc các phương pháp tối ưu cổ điển) để tạo dataset chính sách
``chuyên gia'', sau đó dùng MADDPG fine-tune trên dataset đó để học
chính sách vừa gần nghiệm tìm kiếm vừa khả thi thời gian thực. Hướng
này có tiềm năng thu hẹp khoảng cách giữa heuristic AO-Grid ($5{,}17$
b/s/Hz, một mốc so sánh --- không phải cận trên) và MADDPG hiện tại
($4{,}52$ b/s/Hz, tức $\approx 87\%$ sum-rate của AO-Grid), đồng thời
duy trì độ trễ suy luận cỡ mili-giây.

\textbf{H5. Transformer-based policy cho variable-size $N, K$.} Thay
thế MLP Actor bằng transformer hoặc graph attention network có khả
năng xử lý input có kích thước thay đổi. Điều này giúp một chính sách
được huấn luyện trên một cấu hình có thể tổng quát hóa sang nhiều cấu
hình khác mà không cần huấn luyện lại.

\textbf{H6. Ứng dụng cho ISAC và URLLC.} Kết hợp STAR-RIS RSMA với
cảm biến tích hợp (Integrated Sensing and Communication) hoặc dịch vụ
URLLC mở ra nhiều bài toán mới về cân bằng giữa truyền tin, cảm biến,
và đáp ứng QoS độ trễ thấp \cite{soleymani2024starris},
\cite{adams2024rsma}, \cite{jin2024arisrsma}, \cite{jang2024rsma}.
Trong kịch bản ISAC, hàm phần thưởng cần được mở rộng để bao gồm cả
sensing performance metrics.

\textbf{H7. Federated MADDPG.} Trong các kịch bản multi-BS, mỗi BS có
một MADDPG controller cục bộ. Bài toán đặt ra là làm thế nào để chia
sẻ tri thức học được giữa các BS mà không cần truyền raw data
(privacy-preserving). Federated Learning cung cấp khung tiếp cận tự
nhiên cho vấn đề này \cite{noor2025aimulti}, \cite{mousavi2025fed}.

\textbf{H8. Kiểm chứng phần cứng (Hardware-in-the-loop).} Triển khai
khuôn khổ MADDPG trên SDR testbed (như USRP hoặc GNU Radio) hoặc hệ
thống RIS thực tế để kiểm chứng các kết quả mô phỏng. Đây là bước
thiết yếu trước khi đề xuất đưa vào tiêu chuẩn hóa.

\noindent\hrulefill

\vspace{0.5cm}
Luận văn đã hoàn thành các mục tiêu đặt ra ban đầu và mở ra nhiều hướng
phát triển có giá trị cho cộng đồng nghiên cứu. Tác giả hy vọng các kết
quả của Luận văn sẽ là một bước tiến nhỏ góp phần vào sự phát triển của
mạng vô tuyến thế hệ thứ sáu trong tương lai gần.
